\documentclass[11pt,a4paper]{article}
\usepackage[T1]{fontenc}
\usepackage[utf8]{inputenc}
\usepackage{lmodern}
\usepackage[a4paper,margin=22mm,headheight=15pt,headsep=8mm]{geometry}
\usepackage{amsmath,amssymb,amsthm,mathtools,mathrsfs}
\usepackage{microtype,booktabs,array,longtable,graphicx}
\usepackage[dvipsnames]{xcolor}
\usepackage{enumitem,fancyhdr,xurl}
\usepackage[unicode,colorlinks=true,linkcolor=MidnightBlue,citecolor=MidnightBlue,urlcolor=MidnightBlue]{hyperref}
\hypersetup{pdftitle={Gaussian polylogarithms: proofs, reductions, and exact certificates},
 pdfauthor={Research companion prepared for the ProveIt project},
 pdfsubject={Proofs of Gaussian multiple polylogarithm conjectures and a reproducible source audit}}
\DeclareMathOperator{\Li}{Li}
\renewcommand{\Re}{\operatorname{Re}}
\renewcommand{\Im}{\operatorname{Im}}
\newcommand{\Log}{\operatorname{Log}}
\newcolumntype{P}[1]{>{\raggedright\arraybackslash}p{#1}}
\newtheorem{theorem}{Theorem}[section]
\newtheorem{lemma}[theorem]{Lemma}
\newtheorem{corollary}[theorem]{Corollary}
\newtheorem{proposition}[theorem]{Proposition}
\theoremstyle{definition}
\newtheorem{definition}[theorem]{Definition}
\theoremstyle{remark}
\newtheorem{remark}[theorem]{Remark}
\numberwithin{equation}{section}
\setlength{\emergencystretch}{2em}
\setlength{\parskip}{3pt plus 1pt}
\setlist{itemsep=3pt,topsep=5pt}
\allowdisplaybreaks[2]
\raggedbottom
\pagestyle{fancy}
\fancyhf{}
\fancyhead[L]{\small\sffamily Gaussian polylogarithms: proofs and certificates}
\fancyhead[R]{\small\sffamily ProveIt research companion}
\fancyfoot[C]{\thepage}
\renewcommand{\headrulewidth}{0.3pt}
\title{\vspace{-12mm}\textbf{Gaussian polylogarithms:\\[3pt] proofs, reductions, and exact certificates}\\[7pt]
\large A research continuation and audit of the ProveIt manuscript}
\author{Research companion prepared for the ProveIt project}
\date{9 October 2026\quad\textbullet\quad Version 1.0}

\begin{document}
\maketitle
\thispagestyle{plain}
\begin{abstract}
We prove all five displayed Gaussian weight-six double-polylogarithm
conjectures and all three displayed Gaussian weight-four triple-polylogarithm
conjectures in a fixed revision of the ProveIt manuscript. The double
identities follow from an explicit specialization of the established
polylogarithm parity theorem, with the singular limiting cancellation
justified, and are independently recovered from differential recurrences
and exact double-zeta systems. The same calculation gives a complete
seven-row weight-eight table and rational coefficient formulas at every
even weight. A logarithmic-moment argument gives a closed classical
reduction for every position of a single index $2$ among indices $1$.
It reduces the top two possible series-depth layers at every weight and
explains the signed binomial coefficients of the highest-weight
half-Gaussian value. At the primitive sixth roots, one parity component
is always a rational multiple of the corresponding power of $\pi$.
Four proposed mixed-root directions reduce exactly
to the manuscript's existing single-value constants. We also correct a
color-exchange error in a diagonal antisymmetry claim, a diagonal example,
and a reversed parity comment for shifted polygamma differences.
Finally, we give a rational H\"older certificate with an explicit uniform
tail bound and a polynomial-time exact implementation. The accompanying
package contains source, exact coefficient tables, independent numerical
checks, replayable rational certificates, and a proposed integration patch.
All novelty and independence claims are kept separate from what the
proofs actually establish.
\end{abstract}

\noindent\textbf{Keywords.} Multiple polylogarithms; roots of unity;
Gaussian periods; parity theorem; Nielsen polylogarithms; iterated
integrals; exact rational enclosures.

\medskip
\noindent\textbf{Source baseline.} This article audits the manuscript at
revision
\begin{center}
\small\texttt{09812e81e578c3f13f54b2d0c46a95a0e0665b19}
\end{center}
of \texttt{VladimirReshetnikov/ProveIt}, under
\texttt{Analysis/Polylogarithms/docs/manuscript}. Statements about the
source refer to that revision, not to a moving branch.

\clearpage
\begingroup
\small
\setlength{\parskip}{0pt}
\tableofcontents
\endgroup
\clearpage
\input{sections/introduction}
\input{sections/parity}
\input{sections/one_two}
\input{sections/mixed}
\input{sections/certificates}
\input{sections/corrections}
\input{sections/verification}
\input{sections/research}
\input{references}
\end{document}
